\documentclass[a4paper,10pt]{article}
\usepackage[left=0.5in, right=0.5in, top=0.5in, bottom=0.5in]{geometry}
\usepackage{enumitem}
\usepackage{hyperref}
\usepackage{titlesec}

% Formatação minimalista das seções
\titleformat{\section}{\large\bfseries\uppercase}{}{0em}{}[\titlerule]
\titlespacing{\section}{0pt}{1.5ex}{1ex}

% Formatação dos links
\hypersetup{
    colorlinks=true,
    linkcolor=black,
    filecolor=black,
    urlcolor=black,
}

\begin{document}
\pagestyle{empty} % Remove a numeração das páginas

% -------------------- CABEÇALHO --------------------
\begin{center}
    {\huge \textbf{WESLLEY BATISTA}} \\
    \vspace{4pt}
    \hrule
    \vspace{6pt}
    {\small
    \makebox[\textwidth]{%
        +55 48 98880 5667 \hfill $|$ \hfill
        seawsurf@gmail.com \hfill $|$ \hfill
        Santa Catarina, Brasil \hfill $|$ \hfill
        \href{https://www.linkedin.com/in/weslley-batista-/}{linkedin.com/in/seawsurf} \hfill $|$ \hfill
        \href{https://github.com/sea-surf}{github.com/sea-surf}%
    }
    }
\end{center}

% -------------------- EXPERIÊNCIA --------------------
\section*{Experiência}
\noindent
\textbf{Laboratório de Visão Computacional - Universidade Federal de Santa Catarina} \hfill Santa Catarina, Brasil \\
Estagiário de Pesquisa \hfill Maio de 2026 -- Atual
\begin{itemize}[leftmargin=*, noitemsep, topsep=0pt]
    \item Pesquisou artigos acadêmicos, simuladores de drones no NVIDIA Isaac e conjuntos de dados públicos para treinar um modelo confiável de detecção de fogo e fumaça, usando um conjunto com mais de 100 mil imagens (\url{https://doi.org/10.57760/sciencedb.j00104.00103}) para aplicação em drones.
    \item Apresentou resultados técnicos e atualizações da pesquisa à Global Drones, empresa parceira de pesquisa e desenvolvimento de drones (\url{https://globaldrones.com.br/}).
    \item Projeto de visão computacional no GitHub (\url{https://github.com/sea-surf/computer_vision}).
\end{itemize}


% -------------------- FORMAÇÃO --------------------
\section*{Formação}
\noindent
\textbf{Bacharelado em Engenharia da Computação} \hfill Universidade Federal de Santa Catarina, Brasil \hfill Em andamento\\ \hfill 3º semestre


% -------------------- PERFIL E INTERESSES --------------------
\section*{Perfil e interesses}
\noindent Estudante de Engenharia da Computação. \\
\textbf{Interesses:} Computação gráfica e visão computacional, aprendizado de máquina, ciência de dados, robótica e surfe.

% -------------------- COMPETÊNCIAS --------------------
\section*{Competências}
\begin{itemize}[leftmargin=*, noitemsep, topsep=0pt]
    \item \textbf{Linguagens de programação}: Python, C++
    \item \textbf{Ferramentas e bibliotecas}: Git, visão computacional (YOLO, OpenCV, PyTorch)
\end{itemize}


% -------------------- IDIOMAS --------------------
\section*{Idiomas}
\noindent Inglês, português

\end{document}
